%------------------------------------------------------------------------------
% CV in Latex
% Author : Charles Rambo
% Based off of: https://github.com/sb2nov/resume and Jake's Resume on Overleaf
% Most recently updated version may be found at https://github.com/fizixmastr
% License : MIT
%------------------------------------------------------------------------------

\documentclass[A4,11pt]{article}
%\documentclass[letterpaper,11pt]{article} %For use in US
\usepackage[english, russian]{babel} % для русского языка
\usepackage[utf8]{inputenc}
\usepackage[T2A]{fontenc}
\usepackage{amsmath}
\usepackage{amsthm} %для Теорем
\usepackage{amsmath}
\usepackage{indentfirst} %для отступа в первом абзаце
\usepackage{geometry}
\usepackage{titleps}
\usepackage{graphicx}
\usepackage{soulutf8}
\usepackage{booktabs}
\usepackage{gensymb}
\usepackage{array}
\usepackage{latexsym}
\usepackage[empty]{fullpage}
\usepackage{titlesec}
\usepackage{marvosym}
\usepackage[usenames,dvipsnames]{color}
\usepackage{verbatim}
\usepackage{enumitem}
\usepackage[hidelinks]{hyperref}
\usepackage[english]{babel}
\usepackage{tabularx}
\usepackage{tikz}
\usepackage{fontawesome5}
\usepackage{listings}
\usepackage{xcolor}
\usepackage{enumitem}
\input{glyphtounicode}

\begin{comment}
I am by no means a professional when it comes to the CV's/resumes, I have
received various trainings on how to write a CV and resume from my high
school, as well as the Austin College and University of Eastern Finland's
career counseling departments. As I intend to share my CV as a template, I
feel that it is my responsibility to provide explanations of my work.
\end{comment}


%-----FONT OPTIONS-------------------------------------------------------------
\begin{comment}
The font of the document will impact not just how readable it is, but how it is
perceived. In the "The Craft of Scientific Writing" by Michael Alley, shares a
common fonts for publication as well as their use. I have chosen to use
Palatino for its legibility, some others are given below. There is far too much
about typography to discus here. Note: serif fonts have short projecting
strokes, sans-serif fonts are sans (without) these strokes.
\end{comment}


% serif
 \usepackage{palatino}
% \usepackage{times} %This is the default as well
% \usepackage{charter}

% sans-serif
% \usepackage{helvet}
% \usepackage[sfdefault]{noto-sans}
% \usepackage[default]{sourcesanspro}

%-----PAGE SETUP---------------------------------------------------------------

% Adjust margins
\addtolength{\oddsidemargin}{-1cm}
\addtolength{\evensidemargin}{-1cm}
\addtolength{\textwidth}{2cm}
\addtolength{\topmargin}{-1cm}
\addtolength{\textheight}{2cm}

% Margins for US Letter size
%\addtolength{\oddsidemargin}{-0.5in}
%\addtolength{\evensidemargin}{-0.5in}
%\addtolength{\textwidth}{1in}
%\addtolength{\topmargin}{-.5in}
%\addtolength{\textheight}{1.0in}

\urlstyle{same}

\raggedbottom
\raggedright
\setlength{\tabcolsep}{0cm}

% Sections formatting
\titleformat{\section}{
  \vspace{-4pt}\scshape\raggedright\large
}{}{0em}{}[\color{black}\titlerule \vspace{-5pt}]

% Ensure that .pdf is machine readable/ATS parsable
\pdfgentounicode=1

%-----CUSTOM COMMANDS FOR FORMATTING SECTIONS----------------------------------
\newcommand{\CVItem}[1]{
  \item\small{
    {#1 \vspace{-2pt}}
  }
}

\newcommand{\CVSubheading}[4]{
  \vspace{-2pt}\item
    \begin{tabular*}{0.97\textwidth}[t]{l@{\extracolsep{\fill}}r}
      \textbf{#1} & #2 \\
      \small#3 & \small #4 \\
    \end{tabular*}\vspace{-7pt}
}

\newcommand{\CVSubheadingDenFirst}[2]{
  \vspace{-2pt}\item
    \begin{tabular*}{0.97\textwidth}[t]{l@{\extracolsep{\fill}}r}
      \textbf{#1} & #2
    \end{tabular*}\vspace{-7pt}
}

\newcommand{\CVSubheadingDen}[2]{
  \vspace{-14pt}\item
    \begin{tabular*}{0.97\textwidth}[t]{l@{\extracolsep{\fill}}r}
      \textbf{#1} & #2
    \end{tabular*}\vspace{-7pt}
}

\newcommand{\CVSubSubheading}[2]{
    \item
    \begin{tabular*}{0.97\textwidth}{l@{\extracolsep{\fill}}r}
      \text{\small#1} & \text{\small #2} \\
    \end{tabular*}\vspace{-7pt}
}

\newcommand{\CVSubItem}[1]{\CVItem{#1}\vspace{-4pt}}

\renewcommand\labelitemii{$\vcenter{\hbox{\tiny$\bullet$}}$}

\newcommand{\CVSubHeadingListStart}{\begin{itemize}[leftmargin=0.5cm, label={}]}
% \newcommand{\resumeSubHeadingListStart}{\begin{itemize}[leftmargin=0.15in, label={}]} % Uncomment for US
\newcommand{\CVSubHeadingListEnd}{\end{itemize}}
\newcommand{\CVItemListStart}{\begin{itemize}}
\newcommand{\CVItemListEnd}{\end{itemize}\vspace{-5pt}}

%------------------------------------------------------------------------------
% CV STARTS HERE  %
%------------------------------------------------------------------------------
\begin{document}

%-----HEADING------------------------------------------------------------------
\begin{comment}
In Europe it is common to include a picture of ones self in the CV. Select
which heading appropriate for the document you are creating.
\end{comment}

\begin{minipage}[c]{0.05\textwidth}
\-\
\end{minipage}
\begin{minipage}[c]{0.2\textwidth}
\begin{tikzpicture}
    \node at (0,-.7) {\includegraphics[width = 3cm]{Фото.png}};
    % if necessary the picture may be moved by changing the at (coordinates)
    % width defines the 'zoom' of the picture
\end{tikzpicture}
\hfill\vline\hfill
\end{minipage}
\begin{minipage}[c]{0.5\textwidth}
    \textbf{\Huge \scshape{Мельников Денис}} \\
    % \scshape sets small capital letters, remove if desired
    \small{\faPhone} \hspace{0.1cm} \small{+7 931 511 58 76} \\
    \faEnvelope \hspace{0.1cm} \href{melnikov.denis.aleksandrovich@gmail.com}{\underline{melnikov.denis.aleksandrovich@gmail.com}}\\
    \faTelegram \hspace{0.1cm} \href{https://t.me/daikofnetu} {\underline{@daikofnetu}} \\
    % Be sure to use a professional *personal* email address
    \faGithub \hspace{0.1cm} \href{https://www.github.com/dinichthys}{\underline{Dinichthys}} \\
    % you should adjust you linked in profile name to be professional and recognizable
    \faMapMarker \hspace{0.1cm} {\underline{Москва, Россия}} \\
\end{minipage}

% Without picture
%\begin{center}
%    \textbf{\Huge \scshape Charles Rambo} \\ \vspace{1pt} %\scshape sets small capital letters, remove if desired
%    \small +1 123-456-7890 $|$
%    \href{mailto:you@provider.com}{\underline{you@provider.com}} $|$\\
%    % Be sure to use a professional *personal* email address
%    \href{https://linkedin.com/in/your-name-here}{\underline{linkedin.com/in/charles-rambo}} $|$
%    % you should adjust you linked in profile name to be professional and recognizable
%    \href{https://github.com/fizixmastr}{\underline{github.com/fizixmastr}}
%\end{center}



\begin{comment}
This CV was written for specifically for positions I was applying for in
academia, and then modified to be a template.

A standard CV is about two pages long where as a resume in the US is one page.
sections can be added and removed here with this in mind. In my experience,
education, and applicable work experience and skills are the most import things
to include on a resume. For a CV the Europass CV suggests the categories: Work
Experience, Education and Training, Language Skills, Digital Skills,
Communication and Interpersonal Skills, Conferences and Seminars, Creative Works
Driver's License, Hobbies and Interests, Honors and Awards, Management and
Leadership Skills, Networks and Memberships, Organizational Skills, Projects,
Publications, Recommendations, Social and Political Activities, Volunteering.

Your goal is to convey a who, what , when, where, why for every item you share.
The who is obviously you, but I believe the rest should be done in that order.
For example below. An employer cares most about the degree held and typically
less about the institution or where it is located (This is still good
information though). Whatever order you choose be consistent throughout.
\end{comment}

%-----EDUCATION----------------------------------------------------------------
\section{Образование}
  \CVSubHeadingListStart
%    \CVSubheading % Example
%      {Degree Achieved}{Years of Study}
%      {Institution of Study}{Where it is located}
    \CVSubheading
      {{Бакалавриат $|$ \emph{\small{Второй курс, Средний балл: 8.65}}}}{Сен. 2024 -- Настоящее время}
      {Физтех-школа радиотехники и компьютерных технологий \\ Московский физико-технический институт}{Долгопрудный, Россия}
\CVItemListEnd

\section{Дополнительные курсы}
  \CVSubHeadingListStart
%    \CVSubheading % Example
%      {Degree Achieved}{Years of Study}
%      {Institution of Study}{Where it is located}
    \CVSubheadingDenFirst
      {{Курс Ильи Рудольфовича Дединского по языку C и языкам ассемблера}}{Сен. 2024 -- Май 2025}
    \CVSubheadingDen
      {{Курс Ильи Рудольфовича Дединского по языку C++}}{Сен. 2025 -- Май 2026}
    \CVSubheadingDen
      {{Введение в LLVM IR от Сергея Лисицына (МФТИ)}}{Фев. 2026 -- Март 2026}
    \CVSubheadingDen
      {{Введение в архитектуру вычислительных систем от огранизации СБЕР}}{Фев. 2026 -- Май 2026}
    \CVSubheadingDen
      {{Введение в тензорные компиляторы от огранизации СБЕР}}{Фев. 2026 -- Май 2026}
  \CVSubHeadingListEnd

% %-----WORK EXPERIENCE----------------------------------------------------------
% \begin{comment}
% try to briefly explain what you did and why it is relevant to the position you
% are seeking
% \end{comment}

\section{Опыт работы}
  \CVSubHeadingListStart
%    \CVSubheading %Example
%      {What}{When}
%      {Short Description}{}
    \CVSubheading
      {Опыт работы в команде разработчиков компилятора Go}{Июль - Сентябрь 2025}
      {}{}
      \small{Стажёр (ассистент-инженер) в компании Huawei на полной ставке. \\ Я дорабатывал определённый оптимизационный пасс в компиляторе, чтобы он рассматривал больше случаев. По итогам моей стажировки мой коммит в официальный репозиторий компилятора был принят и изменения внесены.}
  \CVSubHeadingListEnd

\section{Опыт преподавания}
  \CVSubHeadingListStart
%    \CVSubheading %Example
%      {What}{When}
%      {Short Description}{}
    \CVSubheading
      {Опыт менторства на курсе Ильи Рудольфовича Дединского}{Сен. 2025 - Май 2026}
      {}{}
      \small{Я работал ментором на курсе Ильи Рудольфовича Дединского по языку C и языкам ассемблера для первокурсников на протяжении полного года обучения.}
  \CVSubHeadingListEnd

% \section{Work Experience}
%   \CVSubHeadingListStart
% %    \CVSubheading %Example
% %      {What you did}{When you worked there}
% %      {Who you worked for}{Where they are located}
% %      \CVItemListStart
% %        \CVItem{Why it is important to this employer}
% %      \CVItemListEnd
%     \CVSubheading
%       {Integration Engineering Intern}{June 2018 -- August 2019}
%       {Finisar Corp.}{Sherman, TX}
%       \CVItemListStart
%         \CVItem{Worked in ISO 4 cleanroom developing applications to improve efficiency and creating specs}
%         \CVItem{Employed metrology and microscopy for failure analysis and developing process for wet etching}
%         \CVItem{Member of Emergency Response Team}
%       \CVItemListEnd
%     \CVSubheading
%       {Laboratory Assistant}{January 2016 -- July 2016}
%       {North Lake College}{Irving, TX}
%       \CVItemListStart
%         \CVItem{Inventoried and maintained Physics Department lab equipment}
%         \CVItem{Physics tutoring}
%     \CVItemListEnd
%     \CVSubheading
%       {Assistant Manager}{December 2006 -- August 2015}
%       {Sun \& Ski Sports}{Austin, TX}
%       \CVItemListStart
%         \CVItem{Led a team of 20+ employees}
%         \CVItem{Ran social media, as well as all grassroots marketing}
%       \CVItemListEnd
%   \CVSubHeadingListEnd

%-----PROJECTS AND RESEARCH----------------------------------------------------
\begin{comment}
Ideally the title of the work should speak for what it is. However if you feel
like you should explain more about why the project is applicable to this job,
use item list as is shown in the work experience section.
\end{comment}
%  Stack
%  Mandelbrot
%  CrackMe
\section{Проекты}
  \CVSubHeadingListStart
%    \CVSubheading
%      {Title of Work}{When it was done}
%      {Institution you worked with}{unused}
    \CVSubheading
      {{Compiler} $|$ \emph{\small{C, Cmake, NASM}}}{\href{https://github.com/Dinichthys/Compiler}{Compiler on GitHub}}{}{}
      \small{
      \textbf{Цель проекта:} Реализовать компилятор для собственного языка программирования под собственную архитектуру, а также NASM.
Front-end компилятора разбивал текст программы на лексемы, по которым строилось абстрактное синтаксическое дерево (AST) с помощью рекурсивного спуска. Затем оно переводилось в промежуточное представление.

Далее оно оптимизировалось Middle-end'ом, в котором производились следующие оптимизации:
\begin{itemize}[nosep]
  \item \textbf{Dead code elimination:} убирал код, который не был достижим.
  \item \textbf{Common subexpression:} результат одинаковых подвыражений вичислялся единожды.
  \item \textbf{Unused variables:} неиспользуемые переменные удалялись из программы.
\end{itemize}

Затем Back-end транслировал IR в ассемблерные инструкции для функционального симулятора,
который впоследствии преобразовывал их в машинный код и выполнял.
      }
    \CVSubheading
      {{LLVMPass} $|$ \emph{\small{C++, Cmake}}}{\href{https://github.com/Dinichthys/LLVMPass}{LLVMPass on GitHub}}{}{}
      \small{
      \textbf{Цель проекта:}
       Реализовать пасс, который бы строил CFG и DFG, а также инструментировал LLVMIR логирующими инструкциями, а затем собирал полученный профиль и по нему на графе отображал цветами и количеством вызовов горячие и холодные функции и базовые блоки в них.
}
    \CVSubheading
      {{Physics} $|$ \emph{\small{C++, Cmake}}}{\href{https://github.com/Dinichthys/Physics}{Physics on GitHub}}{}{}
      \small{
      \textbf{Цель проекта:} Создать физический конструктор разных объектов с возможностью отражения и преломления в них света, а также возможность добавления плагинов. Таким образом, в конечном варианте было 2 плагина - плагин для графической библиотеки (путём его замены можно было перейти с SFML на SDL, например, не переписывая ничего), а также плагин для дорисовки (функционал для рисования различных фигур поверх экрана).

Реализации моих плагинов можно найти в репозиториях - \href{https://github.com/Dinichthys/DR4Backend}{DR4Backend} (плагин графической библиотеки) и \href{https://github.com/Dinichthys/GeomPrimitiveBackend}{GeomPrimitiveBackend} (плагин дорисовки).
}
    \CVSubheading
      {{SPU} $|$ \emph{\small{C, make}}}{\href{https://github.com/Dinichthys/SPU}{SPU on GitHub}}{}{}
      \small{
      \textbf{Цель проекта:}
       Реализовать функциональный симулятор процессора с тремя компонентами:
\begin{itemize}[nosep]
  \item \textbf{Ассемблер:} Транслирует ассемблерные инструкции в бинарное представление.
  \item \textbf{Дизассемблер:} Восстанавливает ассемблерные инструкции из бинарного представления.
  \item \textbf{Процессор:} Исполняет инструкции, описанные с помощью бинарного представления.
\end{itemize}
%        Implement a functional processor simulator with three components:
% \begin{itemize}[nosep]
%   \item \textbf{Assembler:} Translates assembly code to binary representation.
%   \item \textbf{Disassembler:} Reverts binary representation to assembly code.
%   \item \textbf{Processor:} Executes the binary instructions.
% \end{itemize}
}
    \CVSubheading
      {{Printf} $|$ \emph{\small{Assembly (NASM), make}}}{\href{https://github.com/Dinichthys/Printf}{Printf on GitHub}}{}{}
      \small{
      \textbf{Цель проекта:} Реализовать упрощённую версию функции языка C \textbf{"printf"} на языке ассемблера
      с поддержкой следующих спецификаторов форматной строки:
\lstinline|%c|, \lstinline|%s|, \lstinline|%d|, \lstinline|%b|, \lstinline|%o|, \lstinline|%x|, и \lstinline|%f|.
Определение спецификатора происходило с помощью \textbf{JUMP-таблицы}.\\
% \vspace{0.1cm}
\textbf{Дополнительные реализованные возможности:}
\begin{itemize}[nosep]
    \item \textbf{Вывод в файл:} Поддерживается вывод в файл.
    \item \textbf{Цветной вывод:}
    \begin{itemize}[nosep]
        \item При записи в стандартный поток вывода для смены цвета использовались escape-последовательности.
        \item При записи в файл для задания цвета текста были использованы теги HTML (например, \lstinline|<font color="red">|)
    \end{itemize}
    \item \textbf{Числа с плавающей точкой:} Использовался спецификатор \lstinline|%f| для записи чисел с плавающей точкой.
\end{itemize}
%       Implement a simplified version of the C function \textbf{printf} in assembly, supporting the following format specifiers:
% \lstinline|%c|, \lstinline|%s|, \lstinline|%d|, \lstinline|%b|, \lstinline|%o|, \lstinline|%x|, and \lstinline|%f|.
% The parser uses a \textbf{JUMP-table} to efficiently handle specifier letters.\\
% \vspace{0.2cm}
% \textbf{Additional Features:}
% \begin{itemize}[nosep]
%     \item \textbf{File output:} Writing formatted text to files
%     \item \textbf{Color support:}
%     \begin{itemize}[nosep]
%         \item For console output: Insertion ANSI escape sequences for colored text
%         \item For file output: Usage of HTML-style color tags (e.g., \lstinline|<font color="red">|)
%     \end{itemize}
%     \item \textbf{Floating-point support:} Handles \lstinline|%f| specifier for floating-point numbers
% \end{itemize}

      }
    \CVSubheading
      {{Mandelbrot Set} $|$ \emph{\small{C, CMake}}}{\href{https://github.com/Dinichthys/Mandelbrot}{Mandelbrot Set on GitHub}}{}{}
      \small{
      \textbf{Цель проекта:} Оптимизировать программу отрисовки множества Мандельброта разными способами. \\
% \vspace{0.1cm}
\textbf{Было использовано три метода оптимизации:}
\begin{itemize}[nosep]
    \item \textbf{Массивы:} Использование массивов особенным образом позволило компилятору заменить операции над ними на SIMD-инструкции.
    \item \textbf{SIMD Инструкции:} Использование интринсиков позволило оптимизировать программу, потому что теперь стали исполняться инструкции сразу над набором данных.
    \item \textbf{Раскрутка цикла:} Раскрутка цикла помогла избежать зависимости по данным между инструкциями, из-за чего
    конвейер всегда был загружен.
\end{itemize}
%  Try to optimize the work of the program, which draws Mandelbrot Set and makes calculations for every point on the screen. \\
% \vspace{0.2cm}
% \textbf{There are three methods for optimization:}
% \begin{itemize}[nosep]
%     \item \textbf{Arrays:} Usage of arrays in a specific way, so the compiler translates the operations with it to SIMD instructions.
%     \item \textbf{SIMD Instructions:} Usage of intrinsics will optimize the process because one operation will be used for several data at one moment.
%     \item \textbf{Loop unrolling:} Loop unrolling helps to execute data-independent instructions in a parallel way, because the pipeline of the processor will be full filled of instructions, which won't follow to \textbf{bubble} and \textbf{stall}.
% \end{itemize}
      }
  \CVSubHeadingListEnd

%-----HONORS AND AWARDS--------------------------------------------------------

% \section{Достижения за прошедшие 2 года}
%   \CVSubHeadingListStart
% %    \CVSubheading %Example
% %      {What}{When}
% %      {Short Description}{}
%     \CVSubheading
%       {Открытая олимпиада школьников}{2024}
%       {Победитель математической олимпиады от университета ИТМО.}{}
%     \CVSubheading
%       {Межрегиональная олимпиада школьников имени И.Я. Верченко}{2024}
%       {Призёр олимпиады по математике}{}
%     \CVSubheading
%       {Олимпиада школьников «Физтех»}{2024}
%       {Призёр олимпиады по математике}{}
%     \CVSubheading
%       {Отраслевая физико-математическая олимпиада школьников «Росатом»}{2024}
%       {Призёр олимпиады по математике}{}
%   \CVSubHeadingListEnd

%-----INTERESTS-------------------------------------------------------------------
\begin{comment}
This section is compressed from the various skills sections that Euro CV
recommends.
\end{comment}

\section{Интересы}
 \begin{itemize}[leftmargin=0.5cm, label={}]
    \small{\item{
     {Компиляторы}\\
     {Оптимизации}\\
     {Нейронные сети}\\
     {Операционные системы}
    }}
 \end{itemize}

%------------------------------------------------------------------------------

%-----HARD-SKILLS-------------------------------------------------------------------
\begin{comment}
This section is compressed from the various skills sections that Euro CV
recommends.
\end{comment}

\section{Hard Skills}
 \begin{itemize}[leftmargin=0.5cm, label={}]
    \small{\item{
     \textbf{Языки программирования}{: C, C++, Assembly (NASM, TASM, x86-64), Python, GoLang} \\
     \textbf{Инструменты}{: CMake, Make, IDA, Ghidra, Radare2, edb, gdb}
    }}
 \end{itemize}

%------------------------------------------------------------------------------

%-----LANGUAGE-------------------------------------------------------------------
\begin{comment}
This section is compressed from the various skills sections that Euro CV
recommends.
\end{comment}

\section{Языки}
 \begin{itemize}[leftmargin=0.5cm, label={}]
    \small{\item{
\textbf{Русский:} Родной \\
\textbf{Английский:} B1
    }}
 \end{itemize}

%------------------------------------------------------------------------------

%-----SOFT-SKILLS-------------------------------------------------------------------
\begin{comment}
This section is compressed from the various skills sections that Euro CV
recommends.
\end{comment}

\section{Soft Skills}
 \begin{itemize}[leftmargin=0.5cm, label={}]
    \small{\item{
     {Коммуникабельный}\\
     {Уверенный в себе}\\
     {Ответственный}
    }}
 \end{itemize}

%------------------------------------------------------------------------------

\end{document}

